\subsection{ФСР однородной системы. Теорема о ФСР. Примеры нахождения ФСР}

\begin{definition}
Пусть дана однородная система линейных алгебраических уравнений $AX = \mathbf{0}$, где $n$ --- число неизвестных, а $R = \mathrm{rk}(A)$. Упорядоченный набор векторов-столбцов $\{Y_1, Y_2, \dots, Y_k\}$ называется \textbf{фундаментальной системой решений} (ФСР) этой системы, если:
\begin{enumerate}
    \item Каждый вектор $Y_i$ является решением данной системы.
    \item Система векторов $\{Y_1, Y_2, \dots, Y_k\}$ линейно независима (базис в пространстве решений).
    \item Любое решение системы $X$ может быть представлено в виде их линейной комбинации: $X = c_1 Y_1 + c_2 Y_2 + \dots + c_k Y_k$.
\end{enumerate}
\end{definition}

\begin{theorem}[О количестве векторов в ФСР]
Если ранг матрицы $A$ однородной системы $AX = \mathbf{0}$ равен $R < n$, то любая фундаментальная система решений состоит ровно из $n - R$ решений.
\end{theorem}

\begin{tcolorbox}[title=Доказательство, breakable]
Доказательство состоит из двух частей: существования ФСР из $n-R$ векторов и доказательства того, что количество векторов в любой ФСР неизменно.

\textbf{1. Построение (существование).}
Приведем систему к ступенчатому виду методом Гаусса. Пусть (после возможного переименования) переменные $x_1, \dots, x_R$ --- главные, а $x_{R+1}, \dots, x_n$ --- свободные. Придадим свободным переменным значения, образующие единичную матрицу размера $(n-R) \times (n-R)$:
\[
\begin{pmatrix} x_{R+1} \\ x_{R+2} \\ \dots \\ x_n \end{pmatrix} = 
\begin{pmatrix} 1 \\ 0 \\ \dots \\ 0 \end{pmatrix}, 
\begin{pmatrix} 0 \\ 1 \\ \dots \\ 0 \end{pmatrix}, \dots,
\begin{pmatrix} 0 \\ 0 \\ \dots \\ 1 \end{pmatrix}
\]
Для каждого такого набора однозначно вычисляются значения главных переменных. Полученные векторы $X_1, X_2, \dots, X_{n-R}$ линейно независимы, так как матрица, составленная из них, содержит единичный минор порядка $n-R$ (в строках, соответствующих свободным переменным). Полнота этой системы (то, что любое решение есть их линейная комбинация) следует из однозначности определения главных переменных через свободные в методе Гаусса.

\textbf{2. Единственность количества (через основную лемму о линейной зависимости).}
Пусть $\{X_1, \dots, X_{n-R}\}$ --- ФСР специального вида, построенная выше, а $\{Y_1, \dots, Y_k\}$ --- любая другая ФСР.
\begin{itemize}
    \item Так как $Y_i$ --- линейно независимы и выражаются через $X_j$ (по определению ФСР), то по основной лемме о линейной зависимости $k \le n-R$.
    \item Так как $X_j$ --- линейно независимы и выражаются через $Y_i$ (так как $Y_i$ --- ФСР), то $n-R \le k$.
\end{itemize}
Из двух неравенств получаем $k = n-R$. Теорема доказана.
\end{tcolorbox}

\begin{theorem}[О структуре общего решения]
Общее решение однородной СЛАУ $X_O$ представляется в виде:
\[ X_O = c_1 Y_1 + c_2 Y_2 + \dots + c_{n-R} Y_{n-R} \]
где $\{Y_1, \dots, Y_{n-R}\}$ --- любая ФСР системы, а $c_i$ --- произвольные постоянные.
\end{theorem}

\textbf{Пример нахождения ФСР специального вида:}
Рассмотрим систему:
\[
\begin{cases}
x_1 + x_2 + x_3 = 0 \\
2x_1 - 2x_3 + x_4 = 0
\end{cases}
\]
1. Приведем к ступенчатому виду. Главные переменные: $x_1, x_3$ (первые ненулевые коэффициенты в строках), свободные: $x_2, x_4$.
2. Строим таблицу для нахождения ФСР:
\[
\begin{tabular}{|c|c||c|c|}
\hline
$x_1$ & $x_3$ & $x_2$ & $x_4$ \\
\hline
\hline
-1 & 0 & 1 & 0 \\
\hline
-1/2 & 1/2 & 0 & 1 \\
\hline
\end{tabular}
\]
3. Для первой строки ($x_2=1, x_4=0$): из второго уравнения $-2x_3 + 0 = 0 \implies x_3=0$; из первого $x_1 + 1 + 0 = 0 \implies x_1=-1$. Получаем $Y_1 = (-1, 1, 0, 0)^T$.
4. Для второй строки ($x_2=0, x_4=1$): из второго уравнения $-2x_3 + 1 = 0 \implies x_3=1/2$; из первого $x_1 + 0 + 1/2 = 0 \implies x_1=-1/2$. Получаем $Y_2 = (-1/2, 0, 1/2, 1)^T$.
5. Чтобы убрать дроби, умножим $Y_2$ на 2: $\tilde{Y}_2 = (-1, 0, 1, 2)^T$.
\textbf{Ответ в векторной форме:}
\[ X = c_1 \begin{pmatrix} -1 \\ 1 \\ 0 \\ 0 \end{pmatrix} + c_2 \begin{pmatrix} -1 \\ 0 \\ 1 \\ 2 \end{pmatrix} \]

\begin{tcolorbox}[title=Источники, colback=gray!10]
\begin{itemize}
  \item Лекция №\,3, тема «Фундаментальная система решений (ФСР)», временной отрезок 114:00--124:00
  \item PDF «Линейная алгебра\_compressed (1).pdf», раздел 4 «Однородные СЛАУ», стр. 14--15 (701--702)
  \item PDF «Вопросы», вопрос №\,8
\end{itemize}
\end{tcolorbox}
